\documentclass[10pt]{article}
\usepackage[margin=0.75in]{geometry}
\usepackage{amsmath, amssymb, amsthm}
\usepackage{microtype}
\usepackage{enumitem}
\usepackage{fancyhdr}
\usepackage{titlesec}
\usepackage{tcolorbox}
\usepackage{booktabs}

\linespread{1.08}
\setlength{\parskip}{0.4ex plus 0.2ex minus 0.1ex}

\tcbset{
    boxrule=0.8pt,
    colback=white,
    colframe=black,
    arc=0pt,
    boxsep=3pt,
    left=4pt, right=4pt, top=4pt, bottom=4pt
}

\newtcolorbox{result}{
    boxrule=0pt,
    colback=black!5,
    colframe=white,
    arc=0pt,
    boxsep=2pt,
    left=4pt, right=4pt, top=4pt, bottom=4pt
}

\titleformat{\section}{\large\bfseries}{\thesection.}{0.5em}{}
\titleformat{\subsection}{\normalsize\bfseries}{\thesubsection}{0.5em}{}
\titlespacing*{\section}{0pt}{1.5ex}{1ex}
\titlespacing*{\subsection}{0pt}{1ex}{0.5ex}

\setlist{itemsep=1pt, topsep=3pt, parsep=1pt, leftmargin=1.5em}

\pagestyle{fancy}
\fancyhf{}
\fancyhead[L]{\small EECS 127/227AT Homework 2}
\fancyhead[R]{\small Spring 2026}
\fancyfoot[C]{\small\thepage}
\renewcommand{\headrulewidth}{0.4pt}

\theoremstyle{definition}
\newtheorem{definition}{Definition}
\newtheorem{example}{Example}
\theoremstyle{plain}
\newtheorem{theorem}{Theorem}
\newtheorem{lemma}{Lemma}

\newcommand{\RR}{\mathbb{R}}
\newcommand{\NN}{\mathcal{N}}
\newcommand{\Range}{\mathcal{R}}
\DeclareMathOperator{\rank}{rank}
\DeclareMathOperator{\trace}{tr}
\DeclareMathOperator{\spn}{span}
\newcommand{\vv}[1]{\vec{#1}}

\begin{document}

\begin{center}
\begin{tabular}{lcr}
\textbf{EECS 127/227AT} & Optimization Models in Engineering & UC Berkeley \quad Spring 2026 \\
\textbf{Homework 2} & & \\
\end{tabular}
\end{center}

\begin{tcolorbox}
\textbf{This homework is due at 11 PM on February 6th, 2026.}

\textbf{Submission Format:} Your homework submission should consist of a single PDF file that contains all of your answers (any handwritten answers should be scanned).
\end{tcolorbox}

\section{Proof of the Fundamental Theorem of Linear Algebra}

In this question, we will prove the fundamental theorem of linear algebra. For any $A \in \RR^{m \times n}$, let $\NN(A)$, $\Range(A)$, and $\rank(A)$ denote the null space, range and rank of $A$ respectively. For any subspace, $\mathcal{S}$ with dimension, $\dim(\mathcal{S})$, let $\mathcal{S}^\perp$ denote its the subspace orthogonal to $\mathcal{S}$. The fundamental theorem of linear algebra states that,
\begin{equation}
\NN(A) \oplus \Range\!\left(A^\top\right) = \RR^n.
\end{equation}

The proof technique we employ will first show that,
\begin{equation}
\NN(A) = \Range\!\left(A^\top\right)^\perp.
\end{equation}

Then we will prove that we can find orthonormal vectors $\vv{e}_1, \vv{e}_2, \ldots \vv{e}_n$ such that $\NN(A) = \spn(\vv{e}_1, \vv{e}_2, \ldots, \vv{e}_\ell)$ and $\Range\!\left(A^\top\right) = \spn(\vv{e}_{\ell+1}, \vv{e}_{\ell+2}, \ldots, \vv{e}_n)$. As a corollary we get the rank-nullity theorem:
\begin{equation}
\dim(\NN(A)) + \rank(A) = n.
\end{equation}

\begin{enumerate}[label=(\alph*)]
\item First, show that $\NN(A) \subseteq \Range\!\left(A^\top\right)^\perp$.

\textit{HINT: Consider $\vv{u}$ in $\NN(A)$, $\vv{v} \in \Range\!\left(A^\top\right)$ and show that $\vv{u}^\top \vv{v} = 0$.}

\item Now show that: $\Range\!\left(A^\top\right)^\perp \subseteq \NN(A)$.

\textit{HINT: Show that any vector $\vv{v}$ that is orthogonal to all vectors in the range of $A^\top$ satisfies $A\vv{v} = 0$. To do this, consider $\vv{v} \in R(A^\top)^\perp$ and what it implies for $\vv{v}^\top A^\top$.}

\item Note that we could apply the orthogonal decomposition theorem (theorem 19 in the course reader) at this point to complete the proof. However, instead we'll work through how to re-derive that result directly. Let $\dim(\NN(A)) = \ell$ and let $\vv{e}_1, \ldots, \vv{e}_\ell$ be an orthonormal basis for $\NN(A)$. Consider an extension of the basis to an orthonormal basis, $\vv{e}_1, \ldots, \vv{e}_n$ for $\RR^n$. We will prove that $\vv{e}_{\ell+1}, \ldots, \vv{e}_n$ form a basis for $\Range\!\left(A^\top\right)$ and as a consequence, the dimension of $\Range\!\left(A^\top\right)$ is $n - \ell$.

\begin{enumerate}[label=\roman*.]
\item Show that $\Range\!\left(A^\top\right)$ lies in the span of $\vv{e}_{\ell+1}, \ldots, \vv{e}_n$.

\textit{HINT: Express any vector $\vv{u} \in \Range\!\left(A^\top\right)$ as $\vv{u} = \sum_{i=1}^{n} \alpha_i \vv{e}_i$. What are the values of $\alpha_i$?}

\item From part (i) we know that $\Range\!\left(A^\top\right) \subseteq \spn(\vv{e}_{\ell+1}, \ldots, \vv{e}_n)$, but we want something stronger. Show that in fact $\Range\!\left(A^\top\right) = \spn(\vv{e}_{\ell+1}, \ldots, \vv{e}_n)$.

\textit{HINT: First, prove $\dim\!\left(\Range\!\left(A^\top\right)\right) = \dim(\spn(\vv{e}_{\ell+1}, \ldots, \vv{e}_n)) = n - \ell$ by contradiction. Assume $\dim\!\left(\Range\!\left(A^\top\right)\right) = k < n - \ell$, show that a vector $\vv{u} \in \spn(\vv{e}_{\ell+1}, \ldots, \vv{e}_n)$ and $\vv{u} \notin \Range\!\left(A^\top\right)$ cannot exist.}

\textit{Specifically, let $\vv{f}_1, \vv{f}_2, \ldots, \vv{f}_k$ be an orthonormal basis for $\Range\!\left(A^\top\right)$, we can find non-zero $\vv{u}_\perp = \vv{u} - \sum_{i=1}^{k} (\vv{f}_i^\top \vv{u}) \vv{f}_i$ that is orthogonal to $\Range\!\left(A^\top\right)$. Does $\vv{u}_\perp$ lie in $\NN(A)$? Does $\vv{u}_\perp$ also lie in $\spn(\vv{e}_{\ell+1}, \ldots, \vv{e}_n)$? Does this lead to a contradiction? Think of $n - \ell = 3$ and $k = 2$ for visualization.}

\textit{HINT: Second, you can use without proof the fact that for two subspaces, $\mathcal{S}_1$ and $\mathcal{S}_2$, if $\mathcal{S}_1 \subseteq \mathcal{S}_2$ and $\dim(\mathcal{S}_1) = \dim(\mathcal{S}_2)$ then $\mathcal{S}_1 = \mathcal{S}_2$.}
\end{enumerate}

\item Using part (c) argue why $\NN(A) \oplus \Range\!\left(A^\top\right) = \RR^n$ and why the rank nullity theorem holds.
\end{enumerate}

\section{Eigenvalues of Symmetric Matrices}

Let $A \in \mathbb{S}^n$ (i.e., the set of $n \times n$ real symmetric matrices) with eigenvalues $\lambda_i$. Prove that all of the eigenvalues of $A$ are real (i.e.\ that $\lambda_i \in \RR$ for each $i$).

\textit{HINT: Consider the quantity $(Av)^*v$ for eigenvector $v$ where $*$ denotes the conjugate transpose. Note that this is the Hermitian inner product between $Av$ and $v$.}

\textit{NOTE: This exercise is part of the proof of the spectral theorem.}

\section{Distinct Eigenvalues, Orthogonal Eigenspaces}

Let $A \in \mathbb{S}^n$ (i.e.\ the set of $n \times n$ real symmetric matrices) and $(\lambda_1, \vv{u}_1), (\lambda_2, \vv{u}_2)$, $\lambda_1 \neq \lambda_2$ be distinct eigen-pairs of $A$. Show that $\vv{u}_1^\top \vv{u}_2 = 0$, i.e., eigenspaces corresponding to distinct eigenvalues are mutually orthogonal.

\textit{HINT: First try to prove that $\lambda_1 \vv{u}_1^\top \vv{u}_2 = \lambda_2 \vv{u}_1^\top \vv{u}_2$, then show that this implies $\vv{u}_1^\top \vv{u}_2 = 0$.}

\textit{NOTE: This exercise is part of the proof of the spectral theorem.}

\section{Gram Schmidt}

Any set of $n$ linearly independent vectors in $\RR^n$ could be used as a basis for $\RR^n$. However, certain bases could be more suitable for certain operations than others. For example, an orthonormal basis could facilitate solving linear equations.

\begin{enumerate}[label=(\alph*)]
\item Given a matrix $A \in \RR^{n \times n}$, it could be represented as a multiplication of two matrices
\begin{equation}
A = QR,
\end{equation}
where $Q \in \RR^{n \times n}$ is an orthonormal matrix and $R \in \RR^{n \times n}$ is an upper-triangular matrix. For the matrix $A$, describe how Gram-Schmidt process could be used to find the $Q$ and $R$ matrices, and apply this to
\begin{equation}
A = \begin{bmatrix} 3 & -3 & 1 \\ 4 & -4 & -7 \\ 0 & 3 & 3 \end{bmatrix}
\end{equation}
to find an orthonormal matrix $Q$ and an upper-triangular matrix $R$.

\item Given an invertible matrix $A \in \RR^{n \times n}$ and an observation vector $\vv{b} \in \RR^n$, the solution to the equality
\begin{equation}
A\vv{x} = \vv{b}
\end{equation}
is given as $\vv{x} = A^{-1}\vv{b}$. For the matrix $A = QR$ from part 4(a), assume that we want to solve
\begin{equation}
A\vv{x} = \begin{bmatrix} 8 \\ -6 \\ 3 \end{bmatrix}.
\end{equation}
By using the fact that $Q$ is an orthonormal matrix, find $\vv{v}$ such that
\begin{equation}
R\vv{x} = \vv{v}.
\end{equation}
Then, given the upper-triangular matrix $R$ in part 4(a) and $\vv{v}$, find the elements of $\vv{x}$ sequentially.

\item Given an invertible matrix $B \in \RR^{n \times n}$ and an observation vector $\vv{c} \in \RR^n$, find the computational cost of finding the solution $\vv{z}$ to the equation $B\vv{z} = \vv{c}$ by using the $QR$ decomposition of $B$. Assume that $Q$ and $R$ matrices are available, and adding, multiplying, and dividing scalars take one unit of ``computation''.

As an example, computing the inner product $\vv{a}^\top \vv{b}$ is said to be $O(n)$, since we have $n$ scalar multiplication for each $a_i b_i$. Similarly, matrix vector multiplication is $O(n^2)$, since matrix vector multiplication can be viewed as computing $n$ inner products. The computational cost for inverting a matrix in $\RR^n$ is $O(n^3)$, and consequently, the cost grows rapidly as the set of equations grows in size. This is why the expression $A^{-1}\vv{b}$ is usually not computed by directly inverting the matrix $A$. Instead, the $QR$ decomposition of $A$ is exploited to decrease the computational cost.
\end{enumerate}

\section{Determinants}

Consider a unit box $\mathcal{B}$ in $\RR^2$ --- i.e., the square with corners $\begin{bmatrix} 0 \\ 0 \end{bmatrix}$, $\begin{bmatrix} 0 \\ 1 \end{bmatrix}$, $\begin{bmatrix} 1 \\ 0 \end{bmatrix}$, and $\begin{bmatrix} 1 \\ 1 \end{bmatrix}$. Define $A(\mathcal{B})$ as the parallelogram generated by applying matrix $A$ to every point in $\mathcal{B}$.

\begin{enumerate}[label=(\alph*)]
\item For $A = \begin{bmatrix} 1 & 2 \\ 3 & 4 \end{bmatrix}$, calculate the location of each corner of $A(\mathcal{B})$.

\item Write the area of $A(\mathcal{B})$ as a function of $\det(A)$.

\textit{HINT: How are the basis vectors $\begin{bmatrix} 0 \\ 1 \end{bmatrix}$ and $\begin{bmatrix} 1 \\ 0 \end{bmatrix}$ transformed by the matrix multiplication?}

\item Calculate the area of $A(\mathcal{B})$ for each of the following values of $A$.
\begin{enumerate}[label=\roman*.]
\item $A = \begin{bmatrix} 1 & 2 \\ 3 & 4 \end{bmatrix}$
\item $A = \begin{bmatrix} 2 & 1 \\ 4 & 3 \end{bmatrix}$
\item $A = \begin{bmatrix} 1 & 2 \\ 2 & 4 \end{bmatrix}$
\item $A = \begin{bmatrix} 0 & -1 \\ 1 & 0 \end{bmatrix}$
\end{enumerate}
\end{enumerate}

\section{Homework Process}

With whom did you work on this homework? List the names and SIDs of your group members.

\textit{NOTE: If you didn't work with anyone, you can put ``none'' as your answer.}

\end{document}
